\documentclass[10pt,a4paper]{amsart}
\usepackage{amsmath,amssymb,amsthm,graphicx}
\usepackage{verbatim}
\newenvironment{reference}{\comment}{\endcomment}
%\newenvironment{reference}{}{}
\newenvironment{slogan}{\comment}{\endcomment}
\newenvironment{history}{\comment}{\endcomment}

% For commutative diagrams we use Xy-pic
\usepackage[all]{xy}

% We use 2cell for 2-commutative diagrams.
\xyoption{2cell}
% Theorem environments.
%
\theoremstyle{plain}
\newtheorem{theorem}[subsection]{Theorem}
\newtheorem{proposition}[subsection]{Proposition}
\newtheorem{lemma}[subsection]{Lemma}

\theoremstyle{definition}
\newtheorem{definition}[subsection]{Definition}
\newtheorem{example}[subsection]{Example}
\newtheorem{exercise}[subsection]{Exercise}
\newtheorem{situation}[subsection]{Situation}

\theoremstyle{remark}
\newtheorem{remark}[subsection]{Remark}
\newtheorem{remarks}[subsection]{Remarks}

\numberwithin{equation}{subsection}

% Macros
%
\def\lim{\mathop{\mathrm{lim}}\nolimits}
\def\colim{\mathop{\mathrm{colim}}\nolimits}
\def\Spec{\mathop{\mathrm{Spec}}}
\def\Hom{\mathop{\mathrm{Hom}}\nolimits}
\def\Ext{\mathop{\mathrm{Ext}}\nolimits}
\def\SheafHom{\mathop{\mathcal{H}\!\mathit{om}}\nolimits}
\def\SheafExt{\mathop{\mathcal{E}\!\mathit{xt}}\nolimits}
\def\Sch{\mathit{Sch}}
\def\Mor{\mathop{\mathrm{Mor}}\nolimits}
\def\Ob{\mathop{\mathrm{Ob}}\nolimits}
\def\Sh{\mathop{\mathit{Sh}}\nolimits}
\def\NL{\mathop{N\!L}\nolimits}
\def\CH{\mathop{\mathrm{CH}}\nolimits}
\def\proetale{{pro\text{-}\acute{e}tale}}
\def\etale{{\acute{e}tale}}
\def\QCoh{\mathit{QCoh}}
\def\Ker{\mathop{\mathrm{Ker}}}
\def\Im{\mathop{\mathrm{Im}}}
\def\Coker{\mathop{\mathrm{Coker}}}
\def\Coim{\mathop{\mathrm{Coim}}}

% Boxtimes
%
\DeclareMathSymbol{\boxtimes}{\mathbin}{AMSa}{"02}

%
% Macros for moduli stacks/spaces
%
\def\QCohstack{\mathcal{QC}\!\mathit{oh}}
\def\Cohstack{\mathcal{C}\!\mathit{oh}}
\def\Spacesstack{\mathcal{S}\!\mathit{paces}}
\def\Quotfunctor{\mathrm{Quot}}
\def\Hilbfunctor{\mathrm{Hilb}}
\def\Curvesstack{\mathcal{C}\!\mathit{urves}}
\def\Polarizedstack{\mathcal{P}\!\mathit{olarized}}
\def\Complexesstack{\mathcal{C}\!\mathit{omplexes}}
% \Pic is the operator that assigns to X its picard group, usage \Pic(X)
% \Picardstack_{X/B} denotes the Picard stack of X over B
% \Picardfunctor_{X/B} denotes the Picard functor of X over B
\def\Pic{\mathop{\mathrm{Pic}}\nolimits}
\def\Picardstack{\mathcal{P}\!\mathit{ic}}
\def\Picardfunctor{\mathrm{Pic}}
\def\Deformationcategory{\mathcal{D}\!\mathit{ef}}
\usepackage{cleveref}
\crefname{originallemma}{original lemma}{original lemmas}
\Crefname{originallemma}{Original lemma}{Original lemmas}
\makeatletter
\newlabel{lemma-base-change-fs}{{00TJ}{source}{}{sourcepointer-1}{}}
\newlabel{lemma-base-change-fs@cref}{{[originallemma][0][]00TJ}{[1][0][]source}{}{}{}}
\makeatother
\crefname{originallemma}{original lemma}{original lemmas}
\Crefname{originallemma}{Original lemma}{Original lemmas}
\makeatletter
\newlabel{lemma-characterize-formally-smooth-again}{{031I}{source}{}{sourcepointer-2}{}}
\newlabel{lemma-characterize-formally-smooth-again@cref}{{[originallemma][0][]031I}{[1][0][]source}{}{}{}}
\makeatother
\crefname{originallemma}{original lemma}{original lemmas}
\Crefname{originallemma}{Original lemma}{Original lemmas}
\makeatletter
\newlabel{lemma-differential-seq}{{00RU}{source}{}{sourcepointer-3}{}}
\newlabel{lemma-differential-seq@cref}{{[originallemma][0][]00RU}{[1][0][]source}{}{}{}}
\makeatother
\crefname{originallemma}{original lemma}{original lemmas}
\Crefname{originallemma}{Original lemma}{Original lemmas}
\makeatletter
\newlabel{theorem-ffdescent-projectivity}{{05A9}{source}{}{sourcepointer-4}{}}
\newlabel{theorem-ffdescent-projectivity@cref}{{[originallemma][0][]05A9}{[1][0][]source}{}{}{}}
\makeatother
\title{Formal smoothness descends along a faithfully flat base change}
\author{The Stacks Project Authors}
\date{stacks-project:algebra.tex:local-snapshot-20261012; GNU FDL 1.2 or later}
\begin{document}
\maketitle
\noindent\textit{Editorial note.} All delivered mathematical text is the author's own, copied from the pinned Stacks Project chapter file algebra.tex; no statement or proof step was rewritten, reordered or removed.\par\medskip
\noindent\textit{Editorial note.} Preamble provenance, as required for this source: the frozen snapshot stages the chapter's own preamble file beside it (preamble.tex, 217 lines, 8116 bytes, the same shared Stacks-wide preamble the published chow and etale-cohomology cards use); the first line of the pinned chapter file is \textbackslash{}input\{preamble\}, so no preamble text lives inside the chapter itself. The card takes that file's lines 8-18 (the verbatim package and the three \textbackslash{}comment-based environments reference/slogan/history at lines 9-12, section 13.2, together with the Xy-pic setup) and its lines 154 to the runtime end of the file (the native theorem declarations and mathematical macros). Not one character was changed, no package was added, no macro was defined, and the excerpt serves only as preamble, never as a span coordinate.\par\medskip
\noindent\textit{Editorial note.} The same-chapter pointers are located at their \textbackslash{}label lines inside the pinned chapter file; the cross-chapter pointers are located in the project's own tag list staged in the snapshot (sources/stacks-chow/tags.txt), at the verbatim lines carrying their canonical tags.\par\medskip
\noindent\textit{Editorial note.} The counted proof refers to same-chapter statements by label; they are recorded as navigation stubs pointed at their \textbackslash{}label lines, and the arguments stay inside the counted span.\par\medskip
\noindent\textit{Editorial note.} The counted spans contain no \textbackslash{}cite, so the card carries no bibliography.\par\medskip
\noindent\textit{Allowed standard prerequisites (editorial).}\begin{itemize}
\item Preservation of formal smoothness by base change. Invoked at the start of the counted proof (line 38467-38468). Gives the easy direction of the equivalence.
\item Characterisation of formal smoothness by a lifting condition on a two-term complex. Invoked twice in the counted proof (lines 38480-38500). Turns the descent question into exactness of a sequence after base change.
\end{itemize}\medskip
\ifdefined\hypertarget\hypertarget{sourcepointer-1}{}\fi
\noindent\textit{Source reference.} Original lemma 00TJ: sources/stacks-algebra/pinned-original-source.tex, source lines 38037--38037. Navigation pointer for \textbackslash{}ref\{lemma-base-change-fs\} used inside the counted span (line 38468); the argument that consumes it is retained in the span.\par\medskip
\ifdefined\hypertarget\hypertarget{sourcepointer-2}{}\fi
\noindent\textit{Source reference.} Original lemma 031I: sources/stacks-algebra/pinned-original-source.tex, source lines 38128--38128. Navigation pointer for \textbackslash{}ref\{lemma-characterize-formally-smooth-again\} used inside the counted span (line 38477); the argument that consumes it is retained in the span.\par\medskip
\ifdefined\hypertarget\hypertarget{sourcepointer-3}{}\fi
\noindent\textit{Source reference.} Original lemma 00RU: sources/stacks-algebra/pinned-original-source.tex, source lines 34320--34320. Navigation pointer for \textbackslash{}ref\{lemma-differential-seq\} used inside the counted span (line 38486); the argument that consumes it is retained in the span.\par\medskip
\ifdefined\hypertarget\hypertarget{sourcepointer-4}{}\fi
\noindent\textit{Source reference.} Original lemma 05A9: sources/stacks-algebra/pinned-original-source.tex, source lines 22679--22679. Navigation pointer for \textbackslash{}ref\{theorem-ffdescent-projectivity\} used inside the counted span (line 38495); the argument that consumes it is retained in the span.\par\medskip
\begin{lemma}
\label{lemma-descent-formally-smooth}
Let $R \to S$ be a ring map. Let $R \to R'$ be a faithfully flat ring map.
Set $S' = S \otimes_R R'$. Then $R \to S$ is formally smooth if and only
if $R' \to S'$ is formally smooth.
\end{lemma}

\begin{proof}
If $R \to S$ is formally smooth, then $R' \to S'$ is formally smooth by
Lemma \ref{lemma-base-change-fs}.
To prove the converse, assume $R' \to S'$ is formally smooth.
Note that $N \otimes_R R' = N \otimes_S S'$ for any $S$-module $N$. In
particular $S \to S'$ is faithfully flat also.
Choose a polynomial ring $P = R[\{x_i\}_{i \in I}]$ and a surjection
of $R$-algebras $P \to S$ with kernel $J$. Note that $P' = P \otimes_R R'$
is a polynomial algebra over $R'$. Since $R \to R'$ is flat the kernel
$J'$ of the surjection $P' \to S'$ is $J \otimes_R R'$. Hence the
split exact sequence (see
Lemma \ref{lemma-characterize-formally-smooth-again})
$$
0 \to J'/(J')^2 \to \Omega_{P'/R'} \otimes_{P'} S' \to \Omega_{S'/R'} \to 0
$$
is the base change via $S \to S'$ of the corresponding sequence
$$
J/J^2 \to \Omega_{P/R} \otimes_P S \to \Omega_{S/R} \to 0
$$
see
Lemma \ref{lemma-differential-seq}.
As $S \to S'$ is faithfully flat we conclude two things:
(1) this sequence (without ${}'$) is exact too, and (2)
$\Omega_{S/R}$ is a projective $S$-module. Namely, $\Omega_{S'/R'}$
is projective as a direct summand of the free module
$\Omega_{P'/R'} \otimes_{P'} S'$ and
$\Omega_{S/R} \otimes_S {S'} = \Omega_{S'/R'}$ by what we said above.
Thus (2) follows by descent of projectivity
through faithfully flat ring maps, see
Theorem \ref{theorem-ffdescent-projectivity}.
Hence the sequence
$0 \to J/J^2 \to \Omega_{P/R} \otimes_P S \to \Omega_{S/R} \to 0$
is exact also and we win by applying
Lemma \ref{lemma-characterize-formally-smooth-again}
once more.
\end{proof}

\end{document}
